With Kubernetes emerging as the standard platform for workload orchestration, more and more organizations are moving towards it.
The operator pattern in Kubernetes enables encoding operational knowledge into software.
On the other hand, a popular tool to automate the administration of workloads on legacy virtual machines is Ansible.
It is attempting to solve the problem of infrastructure configuration in order to simplify and scale the management of those systems.
\parencite{poulton2025kubernetes, k8soperators, ansibledocsindex}

However, many organizations are using legacy infrastructure such as virtual machines that cannot be easily containerized or migrated to Kubernetes, leading to a split in their operational model.
With some systems managed through Ansible and some systems managed through Kubernetes, information on how to access which system and how to manage it quickly becomes \gls{tribalknowledge}.
This leads to increased overhead in managing the infrastructure, problems with onboarding new employees and a more complex mental model.
\parencite{platformengineering-architects}

This thesis therefore sets out to unify the management of two systems with opposing philosophies into a single control plane.
For this purpose, a Kubernetes operator is implemented that enables Ansible-managed workloads to be managed through the Kubernetes \gls{api}.
Two central questions will be guiding this thesis:

\begin{itemize}
    \item How can the management of Ansible-managed workloads be integrated into
    Kubernetes?
    \item To what extent can Ansible's imperative imposition model be
    reconciled with Kubernetes' declarative, reconciliation-driven operator
    pattern?
\end{itemize}

To achieve this goal and answer the questions, requirements are defined.
Afterward, based on the requirements, an operator is designed and implemented.
Finally, the implementation is evaluated against the requirements.

Chapters \ref{ch:theory} and \ref{ch:technical} introduce theoretical and technical concepts required for understanding the basis on which this thesis builds.
Chapter \ref{ch:goals} explains the goals of the thesis and the problem this thesis is attempting to solve.
Chapter \ref{ch:analysis} looks at the market and investigates whether there are already similar solutions available which could be adopted.
This is followed by chapter \ref{ch:architecture} detailing the architecture of the operator and the \gls{api} it provides for managing Ansible workloads via the control plane.
Next, chapter \ref{ch:implementation} details the implemented logic and explains how the operator manages workloads.
Finally, chapter \ref{ch:evaluation} provides an evaluation of the implemented logic and discusses whether the goals and the requirements have been fulfilled.
